%! AP Statistics — Unit 1, Lesson 1.1 — Student Packet Cover Sheet
\documentclass[10pt]{article}
\usepackage{apstats-article}
\usepackage{apstats-boxes}

%=============================================================================
\begin{document}

% ── Full-bleed navy banner (cover page uses tikz overlay, not \pageheader) ───
\begin{tikzpicture}[remember picture, overlay]
  \fill[navy] (current page.north west)
    rectangle ([yshift=-0.9in]current page.north east);
\end{tikzpicture}
\vspace{-0.8in}

\begin{center}
  {\color{white}\textbf{\LARGE AP Statistics}}\\[4pt]
  {\color{skymid}\large\textbf{Unit 1: Exploring One-Variable Data}}\\[6pt]
  {\color{white}\textbf{\large Lesson 1.1 \quad Introducing Statistics}}
\end{center}

\vspace{0.22in}

\namedateperiod

\vspace{0.18in}

% ── LEARNING TARGETS ─────────────────────────────────────────────────────────
\begin{learningtargetbox}
\small
\begin{enumerate}[leftmargin=1.5em, itemsep=5pt]
  \item \ldots{} explain what statistics is and describe the 4-step statistical
        problem-solving process.
  \item \ldots{} identify the individual, population, sample, and variable(s)
        in a statistical scenario.
  \item \ldots{} explain why variation in data is the reason statistics is useful,
        and distinguish between a sample and a dataset.
\end{enumerate}
\end{learningtargetbox}

\vspace{0.14in}

% ── TABLE OF CONTENTS ────────────────────────────────────────────────────────
\begin{tocbox}
\small
\renewcommand{\arraystretch}{1.5}
\begin{tabularx}{\linewidth}{@{} c l X r @{}}
\rowcolor{sky}
\textbf{\#} & \textbf{Component} & \textbf{Description} & \textbf{Score} \\
\midrule
1 & Warm-Up       & Prior knowledge + class data collection   & \blank{1.2cm} \\
2 & Guided Notes  & Direct instruction notes \& key concepts   & \blank{1.2cm} \\
3 & Group Activity & Scenario analysis with your partner       & \blank{1.2cm} \\
4 & Exit Ticket   & Individual scenario identification (completeness) & \blank{1.2cm} \\
5 & Homework      & Real-world data identification + reflection & \blank{1.2cm} \\
\midrule
  & \multicolumn{2}{r}{\textbf{Total}} & \blank{1.2cm} \\
\end{tabularx}
\end{tocbox}

\vspace{0.14in}

% ── KEEP IN MIND ─────────────────────────────────────────────────────────────
\begin{remindbox}
\small
These ideas will come up again and again all year. Start building the habit of asking them now.

\vspace{6pt}
\begin{tabularx}{\linewidth}{@{} >{\centering\arraybackslash\columncolor{goldbg}}p{0.06\linewidth}
                                    X @{}}
\textbf{\large!} &
\textbf{Association $\neq$ Causation.}\enspace
A study can show that two things are related without proving one \emph{causes} the other. \\[6pt]
\textbf{\large!} &
\textbf{Sample $\neq$ Data.}\enspace
The sample is the set of \emph{individuals} selected. The data are the \emph{values recorded}
for those individuals. \\[6pt]
\textbf{\large!} &
\textbf{Variation is informative, not a problem.}\enspace
If every individual gave the same answer, we would not need statistics at all. \\[6pt]
\textbf{\large!} &
\textbf{The 4-step process spirals all year.}\enspace
Ask a question $\to$ Collect data $\to$ Analyze $\to$ Interpret. You will use this framework
in every unit. \\
\end{tabularx}
\end{remindbox}

\end{document}
